\setcounter{theorem}{11}
universe (a type and a family over it), $\top$ and $\Sigma$ types,
equality with reflection, and a $\Pi$ type with $\U$-domain.
\pagebreak
\begin{definition}[ToS: the theory of GAT signatures]\label{def:gat}
  \begin{alignat*}{10}
    & \Ty && : \Set                                             && \Sigma && :(A:\Ty)\ra(\Tm\,A\ra\Ty)\ra\Ty \\
    & \Tm && : \Ty\ra\Set                                       && (\blank,\blank) && : (a:\Tm\,A)\times\Tm\,(B\,a)\cong\Tm\,(\Sigma\,A\,B) : \mathsf{fst},\mathsf{snd} \\         
    & \U  && : \Ty                                              && \Pi && : (a:\Tm\,\U)\ra(\Tm\,(\El\,a)\ra\Ty)\ra\Ty \\                                                 
    & \El && : \Tm\,\U\ra\Ty                                    && \lam && : \big((x:\Tm\,(\El\,a))\ra\Tm\,(B\,x)\big)\cong\Tm\,(\Pi\,a\,B) : \blank\cdot\blank \\     
    & \top && : \Ty                                             && \Eq && : (A:\Ty)\ra\,\Tm\,A\ra\,\Tm\,A\ra\Ty \\                                                                 
    & \tt && : \vmathbb{1} \cong \Tm\,\top \hspace{6em} && \mathsf{refl} && : (u = v)\cong\Tm\,(\Eq\,A\,u\,v) : \mathsf{reflect}                                    
  \end{alignat*}
\end{definition}
\setcounter{theorem}{12}
For the theory of SOGAT signatures (ToS$^+$), we add a new universe $\U^+$ of sorts for which
variables are allowed: with the help of these we can write second
order functions. $\U^+$ is a subuniverse of $\U$ (witnessed by
$\el^+$) and has a $\Pi$ type with $\U^+$-domain and $\U$-codomain.
\begin{definition}[ToS$^+$: the theory of SOGAT signatures]\label{def:sogat}
  We extend ToS with the following.
  \begin{alignat*}{10}
    & \U^+ && : \Ty                                   && \pi^+ && : (a^+ : \Tm\,\U^+)\ra\Big(\Tm\,\big(\El\,(\el^+\,a^+)\big)\ra\Tm\,\U\Big)\ra\Tm\,\U \\                                          
    & \el^+ && : \Tm\,\U^+ \ra \Tm\,\U \hspace{1.7em} && \lam^+ && : \big(x:\El\,(\el^+\,a^+)\big)\ra\Tm\,(\El\,(b\,x))\cong\Tm\,\big(\El\,(\pi^+\,a^+\,b)\big) : \blank\cdot^+\blank    
  \end{alignat*}
\end{definition}
The first-order version is \autoref{def:sogat_gat}. A (presentation of
a) GAT is defined as a closed type in the syntax of ToS$^+$.
\setcounter{theorem}{13}
\section{Naive semantics of SOGAT signatures}
\label{sec:naive}

In this section, for any SOGAT signature, we define a notion of
first-order model. The idea is that a model is a category together
with the presheaf interpretation of the signature over that category:
the category of presheaves supports a universe, $\Pi$ types, and so
on, so we directly use these when interpreting the type formers of ToS$^+$.
We assume basic working knowledge of
categories with families (CwFs
[19]).
\begin{definition}[First-order model of ToS]\label{def:gat_gat}
  A first-order model of ToS is a CwF (sorts
  are denoted $\Con$, $\Sub$, $\Ty$, $\Tm$, the empty context is $\diamond$, the empty substitution is $\epsilon:\Sub\,\Gamma\,\diamond$, context extension is
  $\blank\ext\blank : (\Gamma:\Con)\ra\Ty\,\Gamma\ra\Con$ with $\p\circ\blank,\q[\blank] :
  \Sub\,\Delta\,(\Gamma\ext A) \cong
  (\gamma:\Sub\,\Delta\,\Gamma)\times\Tm\,\Delta\,(A[\gamma]) :
  \blank\mathbin{{,}}\blank$) equipped with:
  \begin{itemize}
  \item $\top$ and $\Sigma$ types given by isomorphisms \\
    $
    \tt : \vmathbb{1} \cong \Tm\,\Gamma\,\top, \hspace{4em} (\blank,\blank) : (a:\Tm\,\Gamma\,A)\times\Tm\,\Gamma\,(B[\id,a]) \cong \Tm\,\Gamma\,(\Sigma\,A\,B) : \fst,\snd.
    $
  \item A universe given by $\U : \Ty\,\Gamma$ and $\El : \Tm\,\Gamma\,\U\ra\Ty\,\Gamma$.
  \item A function space with domain in $\U$, that is
    $\Pi:(a:\Tm\,\Gamma\,\U)\ra\Ty\,(\Gamma\ext\El\,a)\ra\Ty\,\Gamma$,
    with an isomorphism $\lam : \Tm\,(\Gamma\ext\El\,a)\,B \cong 
    \Tm\,\Gamma\,(\Pi\,a\,B) : \app$.
  \item A strict equality type $\Eq$ with reflection and uniqueness of
    identity proofs.
  \item All the operations listed above are natural in $\Gamma$.
  \end{itemize}
\end{definition}
\begin{definition}[First-order model of ToS$^+$]\label{def:sogat_gat}
  A first-order model of ToS$^+$ is a
  first-order model of ToS equipped with:
  \begin{itemize}
  \item Another universe $\U^+ : \Ty\,\Gamma$ that is a subuniverse of
    $\U$ i.e.\ $\el^+:\Tm\,\Gamma\,\U^+\ra\Tm\,\Gamma\,\U$.
  \item $\U$ is closed under functions with $\U^+$-domain, i.e.\ 
    $\pi^+ :
    (a^+:\Tm\,\Gamma\,\U^+)\ra\Tm\,\big(\Gamma\ext\El\,(\el^+\,a^+)\big)\,\U\ra\Tm\,\Gamma\,\U$
    with $\lam^+ :
    \Tm\,\big(\Gamma\ext\El\,(\el^+\,a)\big)\,(\El\,b) \cong 
    \Tm\,\Gamma\,\big(\El\,(\pi^+\,a^+\,b)\big) : \app^+$.
  \item All the operations listed above are natural in $\Gamma$.
  \end{itemize}
\end{definition}
\begin{problem}[$\PSh$]\label{prob:presheaf}
  Presheaves over a category $C$ form a CwF equipped with $\top$, $\Sigma$ types,
  an equality type with reflection, $\Pi$ types and a Coquand-universe
  $\U$ with $\code:\Ty\,\Gamma\cong\Tm\,\Gamma\,\U:\El$. Unlike in
  original Definition 11, we omit writing universe indices for
  readibility.
\end{problem}
\begin{proof}[Construction]
  We recall the main parts of the construction
  [31] for fixing notations.  $\Gamma:\Con$ is a
  presheaf, that is a family of sets $\Gamma:\C\ra\Set$ with
  reindexing $\gamma_I[f]_\Gamma : \Gamma\,J$ for $\gamma_I:\Gamma\,I$
  and $f:\C(J,I)$ such that $\gamma_I[f\circ g]_\Gamma =
  \gamma_I[f]_\Gamma[g]_\Gamma$ and $\gamma_I[\id]_\Gamma =
  \gamma_I$. A $\sigma:\Sub\,\Delta\,\Gamma$ is a function $\sigma :
  \Delta\,I\ra\Gamma\,I$ such that $(\sigma\,\delta_I)[f]_\Gamma =
  \sigma\,(\delta_I[f]_\Delta)$. A type $A:\Ty\,\Gamma$ is a dependent
  presheaf containing a family $A : (I:\C)\ra\Gamma\,I\ra\Set$ with
  reindexing $a_I[f]_A : A\,J\,(\gamma_I[f]_\Gamma)$ for
  $a_I:A\,I\,\gamma_I$ and $f:\C(J,I)$ satisfying functoriality. Type
  substitution is $A[\gamma]\,I\,\delta_I :=
  A\,I\,(\gamma\,\delta_I)$. A term $a:\Tm\,\Gamma\,A$ is a function
  $a : (\gamma_I:\Gamma\,I)\ra A\,I\,\gamma_I$ such that
  $(a\,\gamma_I)[f]_A = a\,(\gamma_I[f]_\Gamma)$. Term substitution is
  $a[\gamma]\,\delta_I := a\,(\gamma\,\delta_I)$. The empty context is
  constant unit: $\diamond\,I := \vmathbb{1}$. Context extension is
  pointwise: $(\Gamma\ext A)\,I := (\gamma_I:\Gamma\,I)\times
  A\,I\,\gamma_I$, its universal property is given by projections and
  pairing for metatheoretic $\Sigma$ types. $\top$, $\Sigma$ and $\Eq$
  are pointwise. We have the functor Yoneda $\y : \C\ra\PSh(\C)$
  defined by $\y\,I\,J := \C(J,I)$, and we use this to define the
  universe by $\U\,I\,\gamma_I := \Ty\,(\y\,I)$. We observe that
  $\gamma_I[\blank]_\Gamma : \Sub\,(\y\,I)\,\Gamma$ (forward part of
  Yoneda lemma), and define $\Pi\,A\,B\,I\,\gamma_I := \Tm\,(\y\,I\ext
  A[\gamma_I[\blank]_\Gamma])\,(B[\gamma_I[\blank]_\Gamma\circ\p,\q])$.
\end{proof}
\begin{problem}[Locally representable types]
  The CwF of presheaves can be extended to a CwF$^+$, which means a
  CwF with a subsort of $\Ty$ called $\Ty^+$ and a $\Pi^+$ type with
  domain in $\Ty^+$, i.e.\ $\Pi^+ :
  (A:\Ty^+\,\Gamma)\ra\Ty\,(\Gamma\ext A)\ra\Ty\,\Gamma$ with $\lam^+
  : \Tm\,(\Gamma\ext A)\,B\cong\Tm\,\Gamma\,(\Pi^+\,A\,B) : \app^+$,
  natural in $\Gamma$. $\Ty^+$ is classified by the Coquand universe
  $\U^+$.
\end{problem}
\begin{proof}[Construction]
  An element $A : \Ty^+\,\Gamma$ is an $A : \Ty\,\Gamma$ together with
  $\blank\ext_A\blank : (I:\C)\ra\Gamma\,I\ra\C$ and an isomorphism
  $\p_A\circ\blank,\q_A[\blank]_A : \C(J,I\ext_A\gamma_I) \cong
  (f:\C(J,I))\times A\,J\,(\gamma_I[f]_\Gamma) :
  \blank\mathbin{,_A}\blank$ natural in $J$. So
  $\p_A:\C(I\ext_A\gamma_I,I)$ and
  $\q_A:A\,(I\ext_A\gamma_I)\,(\gamma_I[\p_A]_\Gamma)$. Substitution
  is given by $I\ext_{A[\gamma]}\delta_I := I\ext_A\gamma\,\delta_I$
  and we have $\C(J,I\ext_{A[\gamma]}\delta_I) =
  \C(J,I\ext_A\gamma\,\delta_I) \cong (f:\C(J,I))\times
  A\,J\,(\gamma\,\delta_I[f]_\Gamma) = (f:\C(J,I))\times
  A[\gamma]\,J\,(\delta_I[f]_\Delta)$. We define $\Pi^+$ using the
  $\ext_A$ operator which comes with $A$,
  i.e.\ $\Pi^+\,A\,B\,I\,\gamma_I :=
  B\,(I\ext_A\gamma_I)\,(\gamma_I[\p_A]_\Gamma, \q_A)$,
  $b_{I'}[f]_{\Pi^+\,A\,B} := b_{I'}[f\circ\p_A\mathbin{,_A}\q_A]$,
  $\lam^+\,b\,\gamma_I := b\,(\gamma_I[\p_A]_\Gamma,\q_A)$ and
  $\app^+\,t\,(\gamma_I,a_I) :=
  (t\,\gamma_I)[\id_I\mathbin{,_A}a_I]_B$. Like $\U$,
  $\U^+\,I\,\gamma_I := \Ty^+\,(\y\,I)$.
\end{proof}
\begin{definition}[Naive semantics]\label{def:naive}
  Given a category $\C$, $\PSh(\C)$ is a model of ToS$^+$ choosing $\U := \U$, $\El\,a := \El\,a$, $\Pi\,a\,B :=
  \Pi\,(\El\,a)\,B$, $\U^+ := \U^+$, $\el^+\,a^+ :=
  \code\,(\El^+\,a^+)$, $\pi^+\,a^+\,b :=
  \code\,\big(\Pi^+\,(\El^+\,a^+)\,(\El\,b)\big)$. Recall that a SOGAT
  signature $\Omega$ is an element of $\Ty\,\diamond$ in the syntax of ToS$^+$.
  A naive model of $\Omega$ is a category with a
  terminal object $\diamond$ together with the interpretation of $\Omega$ in
  presheaves over this category, i.e.\ 
  $(\C:\mathsf{Cat}^\diamond)\times\Tm_{\PSh(\C)}\,\diamond\llbracket\Omega
  \rrbracket_{\PSh(\C)}$.
\end{definition}
\setcounter{theorem}{18}
\section{Direct semantics of SOGAT signatures}
\label{sec:warmup}

In this section, we define first-order models of SOGATs using a more
careful version of the presheaf model. We make sure that no
Yoneda-encodings are present in the semantics using the idea of
two-level type theory
[4, 10]
where presheaves over a CwF include a universe of ``inner types''
coming from the CwF. We extend two-level type theory with a separate function space
where the domain is an inner type. This function space is isomorphic
to the usual presheaf function space, but has a simpler semantics.
\begin{problem}[Presheaves over a CwF]\label{prob:psh_gat}
  If $\C$ is a CwF, then $\PSh(\C)$ models ToS without using the usual
  presheaf $\U$ and $\Pi$.
\end{problem}
\begin{proof}[Construction]
  We interpret $\top$, $\Sigma$, $\Eq$ as in \autoref{prob:presheaf},
  but define $\U$, $\El$ and $\Pi$ by $\Ty_\C$, $\Tm_\C$ and
  $\ext_\C$, respectively: $\U\,I\,\gamma_I := \Ty_\C\,I$,
  
  $\El\,a\,I\,\gamma_I := \Tm_\C\,I\,(a\,\gamma_I)$,
  
  $\Pi\,a\,B\,I\,\gamma_I := B\,(I\ext_\C a\,\gamma_I)\,(\gamma_I[\p_\C]_\Gamma,\q_\C)$ with
  $\lam\,b\,\gamma_I := b\,(\gamma_I[\p_\C]_\Gamma,\q_C)$ and
  $\app\,t\,(\gamma_I,a_I) := t\,\gamma_I[\id_I\mathbin{,_\C}a_I]_B$.
\end{proof}
\begin{problem}[Presheaves over a CwF$^+$]
  If the category $\C$ is a CwF$^+$, then the previous model extends
  to a model of ToS$^+$ (\autoref{def:sogat_gat}).
\end{problem}
\begin{proof}[Construction]
  We interpret $\U^+$, $\el^+$ and $\pi^+$ by $\Ty^+_\C$, identity and
  $\Pi^+_\C$, respectively: $\U^+\,I\,\gamma_I := \Ty^+_\C\,I$,
  $\el^+\,a\,\gamma_I := a\,\gamma_I$, $\pi^+\,a\,b\,\gamma_I :=
  \Pi^+\,(a\,\gamma_I)\,\big(b\,(\gamma_I[\p_\C]_\Gamma,\q_\C)\big)$,
  $\lam^+\,t\,\gamma_I :=
  \lam^+_\C\,\big(t\,(\gamma_I[\p_\C]_\Gamma,\q_\C)\big)$,
  $\app^+\,t\,(\gamma_I,a_I) :=
  \app^+_\C\,(t\,\gamma_I)[\id_I\mathbin{,_\C}a_I]_{\Tm_\C}$.
\end{proof}
\begin{definition}[Direct semantics]\label{def:direct}
  A direct model of a SOGAT signature $\Omega$ is a category $\C$ with a
  terminal object together with the interpretation of $\Omega$ in
  presheaves over presheaves over $\C$, evaluated at the terminal
  presheaf: $(\C:\mathsf{Cat}^\diamond) \times
  \llbracket\Omega\rrbracket_{\PSh(\PSh(\C))}\,\diamond_{\PSh(\C)}\,{\star}$. Note
  that this makes sense because
  $\PSh(\C):\mathsf{CwF}^+$, hence $\PSh\big(\PSh(\C)\big)$ is a model
  of ToS$^+$.
\end{definition}
\setcounter{theorem}{21}
\begin{theorem}\label{thm:naive_direct}
  For any signature, the naive and direct semantics result in
  isomorphic notions of models.
\end{theorem}
\setcounter{theorem}{21}
\begin{proof}
  We fix a $\C:\mathsf{Cat}^\diamond$, and denote $\D := \PSh(\C)$ and
  $\E := \PSh(\D)$. $\D$ is a model of ToS$^+$
  via \autoref{def:naive} and $\E$ is a model via
  \autoref{def:direct}, and Yoneda navigates between them (it is not
  only a functor, but a CwF pseudomorphism
  [37]). By induction on the syntax of ToS$^+$, we define $\alpha$ for contexts,
  substitutions, types and terms: 
  $\alpha_\Gamma : \Sub_{\E}\,\llbracket\Gamma\rrbracket_\E\,(\y\,\llbracket\Gamma\rrbracket_\D)$,
  $\alpha_\gamma : \alpha_\Gamma\circ\llbracket\gamma\rrbracket_\E = \y\,\llbracket\gamma\rrbracket_\D\circ\alpha_\Delta$,
  $\alpha_A : \llbracket A\rrbracket_\E \cong \y\,\llbracket A\rrbracket_\D[\alpha_\Gamma]$,
  $\alpha_a : \alpha_A[\id,\llbracket a\rrbracket_\E] = \y\,\llbracket a\rrbracket_\D[\alpha_\Gamma]$.
  




  For a signature $\Omega:\Ty\,\diamond$, we thus obtain $\llbracket\Omega\rrbracket_\E\,\diamond_\D\,{\star} \cong \y\,\llbracket\Omega\rrbracket_\D[\alpha_\Gamma]\,\diamond_\D\,{\star} = \Tm_\D\,\diamond_\D\,\llbracket\Omega\rrbracket_\D$.
\end{proof}
Note that there is no size issue when stating the isomorphism because
even if $\E$ is one level up compared to $\D$, we only use small
components from $\E$ when evaluating into it.
